\documentclass[11pt]{amsart}
\usepackage[margin=1.2in]{geometry}
\usepackage{amsmath,amssymb,amsthm}
\usepackage[colorlinks=true,linkcolor=blue!60!black,citecolor=blue!60!black,urlcolor=blue!60!black]{hyperref}
\usepackage{xcolor}
\usepackage{booktabs}

\newtheorem{theorem}{Theorem}[section]
\newtheorem{proposition}[theorem]{Proposition}
\newtheorem{lemma}[theorem]{Lemma}
\newtheorem{corollary}[theorem]{Corollary}
\theoremstyle{definition}
\newtheorem{construction}[theorem]{Construction}
\newtheorem{convention}[theorem]{Convention}
\theoremstyle{remark}
\newtheorem{remark}[theorem]{Remark}
\newtheorem{warning}[theorem]{Warning}

\newcommand{\PP}{\mathbf{P}}
\newcommand{\CC}{\mathbf{C}}
\newcommand{\ZZ}{\mathbf{Z}}
\newcommand{\QQ}{\mathbf{Q}}
\newcommand{\Res}{\operatorname{Res}}
\newcommand{\adj}{\operatorname{adj}}
\newcommand{\Pic}{\operatorname{Pic}}
\newcommand{\Br}{\operatorname{Br}}
\newcommand{\divi}{\operatorname{div}}
\newcommand{\Hx}{H_x}
\newcommand{\Hy}{H_y}
\newcommand{\calR}{\mathcal{R}}
\newcommand{\repo}{\texttt{github.com/mattrobball/unirational}}

\title[Unirationality of general $(2,3)$ hypersurfaces in $\PP^2\times\PP^2$]
{Unirationality of the general $(2,3)$ hypersurface in $\PP^2\times\PP^2$
by tangent residuals:\\ the construction and its verification artifacts}
\date{July 22, 2026}

\begin{document}

\begin{abstract}
Let $X\subset\PP^2_x\times\PP^2_y$ be a general hypersurface of bidegree
$(2,3)$ over $\CC$.  Projection to the second factor makes $X$ a standard
conic bundle over $\PP^2$ with smooth discriminant curve of degree $9$ ---
the smallest discriminant degree for which unirationality of conic bundles
over $\PP^2$ was not previously known, and a case that published heuristics
expected to be non-unirational.  We give a complete account of an explicit
construction proving that $X$ is unirational of degree $20$: the vertical
surface $S=X\cap(\PP^2\times L)$, rational by Tsen's theorem, is mapped
birationally onto a \emph{horizontal} divisor $T\sim 10\Hx+\Hy$ by taking
fiberwise tangent residuals in the second fibration of $X$, whose fibers are
plane cubics.  The key algebraic input is a universal resultant identity
producing a covariant line through the three tangent residuals.  Combined
with known results, the construction yields smooth threefolds that are
unirational, irrational, and not stably rational, with trivial Brauer group
and trivial degree-three unramified cohomology.  These notes are
self-contained on the mathematics and additionally document every machine
computation the proof relies on, with pointers to the certificate files in
the accompanying repository.
\end{abstract}

\maketitle
\tableofcontents

\section{Introduction and statement}\label{sec:intro}

A \emph{standard conic bundle} over $\PP^2$ is a smooth projective threefold
$X$ with a flat morphism $\pi\colon X\to\PP^2$ whose fibers are plane
conics, with relative Picard rank one.  Its discriminant
$\Delta\subset\PP^2$ is a curve; the birational geometry of $X$ is largely
governed by $\deg\Delta$.  Rationality is classical: $X$ is rational iff
$\deg\Delta\le 4$, or $\deg\Delta=5$ with even theta-characteristic
(Beauville, Shokurov; see \cite{Prokhorov} for a survey), and irrational for
smooth $\Delta$ of degree $\ge 6$ by the intermediate Jacobian/Prym
criterion \cite[Th\'eor\`eme~4.9]{Beauville}.  \emph{Unirationality} is far
more delicate.  The published record reaches discriminant degree $8$:
conic bundles over $\PP^2$ with $\deg\Delta\le 8$ are unirational
(Koll\'ar--Mella \cite{KollarMella} together with
\cite[\S14.3]{Prokhorov}; also the preprint \cite{Mella}, which further
covers degree $9$ with \emph{singular} discriminant).  For large general
discriminants the expectation in the literature was negative:
``\emph{It is likely that a conic bundle with sufficiently large and
general discriminant curve is not unirational}''
\cite[\S14.3]{Prokhorov}, echoed in \cite{MassarentiMella}.

The general hypersurface $X$ of bidegree $(2,3)$ in $\PP^2\times\PP^2$ is a
standard conic bundle over the second factor with \emph{smooth} discriminant
of degree $9$: precisely the first case beyond the published boundary.
These notes give a complete, explicit proof of:

\begin{theorem}\label{thm:main}
The general hypersurface $X\subset\PP^2_x\times\PP^2_y$ of bidegree $(2,3)$
over $\CC$ is unirational: there is a dominant rational map
$\PP^3\dashrightarrow X$ of degree $20$, produced by an explicit
construction (Construction~\ref{con:map}).
\end{theorem}

\begin{corollary}\label{cor:invisible}
The very general $(2,3)$ hypersurface $X$ is simultaneously unirational
{\rm(Theorem~\ref{thm:main})}, irrational
{\rm\cite[Th\'eor\`eme~4.9]{Beauville}}, and not stably rational
{\rm\cite[Theorem~1.1]{BvB}}.  Moreover $\Br(X)=0$ and
$H^3_{\mathrm{nr}}(X,\QQ/\ZZ(2))=0$ {\rm\cite{CTV}}, so the failure of
stable rationality is invisible to the classical unramified invariants;
unlike the Artin--Mumford threefold {\rm\cite{ArtinMumford}}, whose
non-stable-rationality is detected by $\Br$, these examples are detected
only by degeneration of the diagonal {\rm\cite{BvB}}.
\end{corollary}

The construction is elementary in style and could have been carried out
classically; the nearest antecedent is the 1980 note of Del Centina--Verdi
\cite{DCV}, which asserted a \emph{special} smooth unirational $(3,2)$
example by an entirely different method (rationality of one double plane
with special branch sextic).  Their displayed data contain
inconsistencies; a corrected pencil is certified in the repository
(\S\ref{sec:artifacts}).  The fiberwise tangent-residual mechanism used
here appears to be new.

\subsection*{Provenance of these notes}
The construction and proof below were produced and mechanically verified in
a computational project archived at \repo.  The mathematical exposition is
self-contained; \S\ref{sec:artifacts} is a complete ledger of the machine
artifacts the proof relies on, what each verifies, and which steps rest on
standard prose mathematics instead.

\section{Setup}\label{sec:setup}

\begin{convention}
We work over $\CC$ (any algebraically closed field of characteristic $0$
works verbatim).  Coordinates $x=(x_0{:}x_1{:}x_2)$ on $\PP^2_x$ and
$y=(U{:}V{:}W)$ on $\PP^2_y$.  A bidegree-$(2,3)$ form is written
\[
F(x,y)=x^{\mathsf T}A(y)\,x,
\]
with $A(y)$ a symmetric $3\times3$ matrix of cubic forms in $y$.  We write
$X=\{F=0\}$ and let $\Hx,\Hy$ denote the restrictions of the two hyperplane
classes.  ``General'' always means: in a dense Zariski-open subset of the
parameter space $\PP^{59}=\PP H^0(\mathcal O_{\PP^2\times\PP^2}(2,3))$,
intersected with the smooth locus.
\end{convention}

The threefold $X$ carries \emph{two} fibrations:
\begin{itemize}
\item $\pi=\mathrm{pr}_y\colon X\to\PP^2_y$, a conic bundle with
discriminant $\Delta=\{\det A(y)=0\}$ of degree $9$, smooth for general
$F$;
\item $\rho=\mathrm{pr}_x\colon X\to\PP^2_x$, a fibration whose fiber over
$x$ is the plane cubic $C_x=\{F(x,-)=0\}\subset\PP^2_y$, smooth for
general $x$.
\end{itemize}
The interplay of the two is the whole story.

\begin{lemma}\label{lem:pic}
For smooth $X$: $\Pic(X)=\ZZ \Hx\oplus\ZZ \Hy$ and
$K_X=\mathcal O_X(-1,0)$.  Every effective divisor class $a\Hx+b\Hy$ has
$a,b\ge0$.  Writing $f_y$ (a conic fiber of $\pi$) and $c_x$ (a cubic fiber
of $\rho$) for the two curve classes,
\[
\Hx\cdot f_y=2,\quad \Hy\cdot f_y=0,\qquad
\Hx\cdot c_x=0,\quad \Hy\cdot c_x=3 .
\]
Consequently a surface $S\in|a\Hx+b\Hy|$ has degree $2a$ over $\PP^2_y$ and
degree $3b$ over $\PP^2_x$.  In particular every multisection of the conic
bundle has even degree, and $\rho$ has no rational section.
\end{lemma}

\begin{proof}
Grothendieck--Lefschetz for the ample divisor $X$ in $\PP^2\times\PP^2$
($\dim X=3$) gives $\Pic$; adjunction gives $K_X$.  Effectivity: from
$0\to\mathcal O(a-2,b-3)\to\mathcal O(a,b)\to\mathcal O_X(a,b)\to0$ and
$H^1(\PP^2\times\PP^2,\mathcal O(m,n))=0$ for all $m,n$, one gets
$h^0(\mathcal O_X(a,b))=h^0(\mathcal O(a,b))-h^0(\mathcal O(a-2,b-3))$,
nonzero only if $a,b\ge0$.  The intersection numbers are immediate.  The
parity statement is consistent with, and geometrically explains, the
nontriviality of the $2$-torsion Brauer class of the generic conic
(restriction--corestriction).
\end{proof}

\begin{remark}
Since every surface in the threefold $X$ is a divisor,
Lemma~\ref{lem:pic} constrains all of them.  Adjunction gives
$K_S=((a-1)\Hx+b\Hy)|_S$ for $S\in|a\Hx+b\Hy|$; for a \emph{horizontal}
class ($a\ge1$) this is effective or trivial, so \emph{no smooth member of
any horizontal linear system is rational}.  Any horizontal surface with
rational normalization --- which Theorem~\ref{thm:main} requires, by
Proposition~\ref{prop:engine} below --- is necessarily a singular member of
its system.  This is the precise sense in which the problem resists naive
searches, and why the construction below must produce a highly singular
divisor.
\end{remark}

\begin{proposition}[Multisection principle]\label{prop:engine}
Let $T\subset X$ be an irreducible surface dominating $\PP^2_y$ with
rational normalization $\widetilde T$, of degree $N$ over $\PP^2_y$.  Then
$X$ is unirational of degree $N$.  Conversely, if $X$ is unirational it
contains such surfaces.
\end{proposition}

\begin{proof}
Base-change: $X\times_{\PP^2_y}\widetilde T\to\widetilde T$ is a conic
bundle over a rational surface carrying the tautological section
$t\mapsto(t,t)$.  A conic with a rational point over the field
$\CC(\widetilde T)$ is isomorphic to $\PP^1$, so
$X\times_{\PP^2_y}\widetilde T$ is birational to
$\widetilde T\times\PP^1$, a rational threefold; the projection to $X$ is
dominant of degree $N$.  For the converse, the image of a general plane
under a dominant map $\PP^3\dashrightarrow X$ is a horizontal surface with
rational (indeed unirational, hence by Castelnuovo rational)
normalization.
\end{proof}

Harris's formulation of the problem
--- find rational surfaces in $X$ \emph{other than} the preimages of
rational curves in $\PP^2_y$ \cite[p.~24]{Harris} --- asks exactly for the
horizontality in Proposition~\ref{prop:engine}: the preimage
$\pi^{-1}(\text{curve})$ is vertical and useless.  The point of the
construction below is that one of these ``useless'' vertical surfaces can
be \emph{converted} into a horizontal one using the second fibration
$\rho$.

\section{The construction}\label{sec:construction}

\begin{construction}\label{con:map}
Fix a general line $L\subset\PP^2_y$ and choose coordinates with
$L=\{W=0\}$.
\begin{enumerate}
\item[\textbf{(I)}] \textbf{The vertical seed.}  Let
$S=X\cap(\PP^2_x\times L)\in|\mathcal O_X(0,1)|$.  Projection to
$L\simeq\PP^1$ makes $S$ a conic bundle over $\PP^1$; its generic fiber is
a smooth conic over the $C_1$ field $\CC(L)\simeq\CC(t)$, hence has a
rational point by Tsen's theorem, and $S$ is a \emph{rational} surface
(irreducible, since the generic fiber is geometrically integral).  Under
the other projection, $S$ is a \emph{trisection} of $\rho$: for general
$x$, the fiber is $C_x\cap L=\{p_1,p_2,p_3\}$.
\item[\textbf{(II)}] \textbf{Tangent residuals.}  For general $x$ the
cubic $C_x$ is smooth at each $p_i$ (this needs only that the binary cubic
$g_x=F(x,-)|_L$ has distinct roots).  Define
\[
\tau\colon S\dashrightarrow X,\qquad
(x,p)\longmapsto\bigl(x,\ \tau_x(p)\bigr),\qquad
T_pC_x\cdot C_x=2p+\tau_x(p).
\]
In the group law of $C_x$ (any flex as origin) this is $p\mapsto -2p$.
\item[\textbf{(III)}] \textbf{The horizontal image.}  Let
$T=\overline{\tau(S)}$.  Theorem~\ref{thm:class} below shows
$T\in|10\Hx+\Hy|$ and $\tau\colon S\dashrightarrow T$ is birational.  By
Lemma~\ref{lem:pic}, $T$ has degree $T\cdot f_y=20$ over $\PP^2_y$; it is
horizontal with rational normalization, so Proposition~\ref{prop:engine}
gives a dominant degree-$20$ map
$\widetilde T\times\PP^1\dashrightarrow X$, proving
Theorem~\ref{thm:main}.
\end{enumerate}
\end{construction}

Everything therefore rests on identifying the image divisor exactly.  This
is done by a universal polynomial identity.

\subsection{The universal resultant identity}

Write a ternary cubic in the coordinates adapted to $L=\{W=0\}$ as
\begin{equation}\label{eq:f}
f(U,V,W)=g(U,V)+W\,h(U,V)+W^2(iU+jV)+kW^3,
\end{equation}
with $g$ a binary cubic and $h$ a binary quadratic.  For
$p=[u{:}v{:}0]$ with $g(u,v)=0$, the tangent line to $\{f=0\}$ at $p$,
evaluated at $[U{:}V{:}W]$, is
\[
P_{U,V,W}(u,v)=U\,g_u(u,v)+V\,g_v(u,v)+W\,h(u,v),
\]
because $f_U,f_V,f_W$ restrict on $W=0$ to $g_u,g_v,h$.  Set
\[
\calR_f(U,V,W)=\Res_{u,v}\bigl(g,\;P_{U,V,W}\bigr),
\]
the resultant of binary forms of degrees $3$ and $2$; up to the standard
normalization, $\calR_f$ is the product of the three tangent lines at the
points of $\{g=0\}$, a cubic in $(U,V,W)$.

\begin{theorem}[Universal identity]\label{thm:identity}
Let $\Delta(g)$ denote the classical discriminant of the binary cubic
$(18abcd-4b^3d+b^2c^2-4ac^3-27a^2d^2$ for $g=au^3+bu^2v+cuv^2+dv^3)$.
Then, identically in the ten coefficients of $f$,
\begin{equation}\label{eq:identity}
\calR_f+\Delta(g)\,f=W^2\,q_f,
\end{equation}
where $q_f$ is a \emph{linear} form in $(U,V,W)$ whose three coefficients
are homogeneous of degree $5$ in the coefficients of $f$ and have greatest
common divisor $1$ in the universal coefficient ring.
\end{theorem}

\begin{proof}
Expand the left side in $W$.  The $W^0$ term is
$\Res(g,\,Ug_u+Vg_v)+\Delta(g)g(U,V)$.  By Euler's relation
$ug_u+vg_v=3g$, each factor $Ug_u(p_i)+Vg_v(p_i)$ of the resultant is a
linear form vanishing at the root $p_i$, so the product is proportional to
$g(U,V)$; the constant is $-\Delta(g)$ (a classical covariant identity,
verified symbolically in the artifacts).  The $W^1$ coefficient of
$\calR_f$ is a binary quadratic in $(U,V)$; evaluating at a root $p_k$ of
$g$ kills every term of the root-expansion except the one attached to
$p_k$, and comparison at the three roots identifies it with
$-\Delta(g)h(U,V)$ (two binary quadratics agreeing at three distinct
points are equal; the case $\Delta(g)=0$ follows by continuity of a
polynomial identity).  Hence the left side of \eqref{eq:identity} is
divisible by $W^2$, and being cubic, the quotient $q_f$ is linear.  The
coefficient degrees: $\Res_{3,2}$ has bidegree $(2,3)$ in the coefficients
of its two arguments, and the coefficients of $P$ are linear in those of
$f$, giving degree $2+3=5$; likewise $\deg\Delta(g)\cdot f=4+1$.
Primitivity ($\gcd=1$) is a machine fact
(\S\ref{sec:artifacts}, artifact~A1/A4).
\end{proof}

\begin{remark}[Geometric content]\label{rem:geom}
Restricted to the cubic $C=\{f=0\}$, the product of the three tangent
lines has divisor $\sum_i(2p_i+r_i)$ where $r_i=\tau(p_i)$, while $W$ has
divisor $p_1+p_2+p_3$.  On the cubic, $2p_i+r_i\sim \Hy|_C$ and
$\sum p_i\sim \Hy|_C$, whence
\[
r_1+r_2+r_3\;\sim\;3\Hy|_C-2\Hy|_C=\Hy|_C .
\]
A degree-$3$ effective divisor on a smooth plane cubic lying in $|\Hy|_C|$
is a line section: \emph{the three tangent residuals are collinear}, and
$\{q_f=0\}$ is the line through them.  (Group-law formulation: the
residual map is $p\mapsto-2p$ and
$(-2p_1)+(-2p_2)+(-2p_3)=-2(p_1+p_2+p_3)=0$.)
\end{remark}

\subsection{The image divisor and its class}

Now let the coefficients of $f=F(x,-)$ vary: they are \emph{quadratic}
forms in $x$.  Then $\Delta(g_x)$ has bidegree $(8,0)$, $\calR$ has
bidegree $(10,3)$, and the covariant line $q=q_{F(x,-)}(y)$ has bidegree
$(10,1)$ --- linear in $y$, with three coefficient forms
$q_U,q_V,q_W\in\CC[x]_{10}$.

\begin{theorem}\label{thm:class}
For general $F$:
\begin{enumerate}
\item On $X$ the identity \eqref{eq:identity} reads $\calR=W^2q$, giving
the exact divisor equality
\[
\divi_X(\calR)=2S+T_q,\qquad
T_q:=\divi_X(q)\in|10\Hx+\Hy|,\qquad (10,3)=2\,(0,1)+(10,1).
\]
\item $T_q$ has no component vertical over $\PP^2_x$; its fiber over a
general $x$ is exactly the three distinct tangent residuals
$\{r_1,r_2,r_3\}$, reduced.
\item $T_q=T=\overline{\tau(S)}$ is irreducible and
$\tau\colon S\dashrightarrow T$ is birational.
\end{enumerate}
\end{theorem}

\begin{proof}
(1) $q|_X\not\equiv0$ since bidegree $(10,1)$ admits no factor of bidegree
$(2,3)$, and then $\calR|_X=W^2q|_X\not\equiv0$; $W$ cuts $S$ (reduced and
irreducible for general $F,L$).

(2) A component of $T_q$ vertical over $\PP^2_x$ (i.e.\ of class $(a,0)$)
would force $q(x,-)\equiv0$ along a curve of $x$'s, i.e.\ a common factor
of $q_U,q_V,q_W$.  For each potential common-factor degree $e\ge1$, the
locus of coefficient triples with such a factor is the image of the proper
multiplication map
$\PP(\CC[x]_e)\times\PP(\CC[x]_{10-e})^{?}\to\PP(\CC[x]_{10})^{3}$-incidence,
hence closed; there are finitely many $e$; so
$\{F:\gcd(q_U,q_V,q_W)=1\}$ is open.  It is nonempty by an explicit
witness (artifact A1/A4), hence dense.  Then $q(x,-)\equiv 0$ for only
finitely many $x$, and no vertical divisorial component exists.  For the
fiber: over general $x$, $C_x\cap\{q(x,-)=0\}$ is three points by
B\'ezout, and by Remark~\ref{rem:geom} the three residuals $r_i$ lie
there; that they are pairwise distinct and off $L$ for general $x$ is an
open condition on $F$ satisfied at an explicit witness (see
Lemma~\ref{lem:openness} and artifact A1/A4), so the fiber is exactly
$\{r_1,r_2,r_3\}$, reduced.

(3) $\overline{\tau(S)}\subseteq T_q$ by \eqref{eq:identity}; both have
degree $3$ over $\PP^2_x$ ($T_q\cdot c_x=3$ by Lemma~\ref{lem:pic}), and
by (2) the generic fiber of $T_q$ is exactly the injective image of the
generic fiber of $S$.  Hence $T_q$ is the closure of the image of the
irreducible surface $S$ (irreducible), and $\tau$ is generically injective
of degree $1$.
\end{proof}

\begin{lemma}[Openness by properness]\label{lem:openness}
Let $\mathcal U\subset\PP^{59}\times\PP^2_x$ be the open locus of pairs
$(F,x)$ where $g_x$ has distinct roots, $q_{F(x,-)}\ne0$, the three
residuals are pairwise distinct, and none lies on $L$.  Then
$\{F:\exists x,\ (F,x)\in\mathcal U\}$ is open in $\PP^{59}$: it is the
complement of the image of the closed set
$(\PP^{59}\times\PP^2)\smallsetminus\mathcal U$ under the proper
projection to $\PP^{59}$.  A single explicit pair $(F_0,x_0)\in\mathcal U$
therefore certifies the condition for general $F$.  Concretely, all four
conditions are checked at a witness by: $\operatorname{disc}(g_{x_0})\ne0$;
$q\neq0$; $\operatorname{disc}\bigl(F(x_0,-)|_{\{q=0\}}\bigr)\ne0$ (three
distinct residuals, no root extraction needed); and
$F(x_0,-)\ne0$ at the point $\{q=0\}\cap L$.
\qed
\end{lemma}

\begin{remark}[Independent Chow-theoretic check of the class]
In $\PP^2_x\times L_p\times\PP^2_r$ with hyperplane classes $H,P,R$, the
incidence $\{p\in C_x\cap L,\ r\in T_pC_x\cap C_x\}$ is a complete
intersection of classes $(2H{+}3P)$, $(2H{+}2P{+}R)$, $(2H{+}3R)$; pushing
down $L_p$ gives $20H^2+36HR+9R^2$.  Subtracting twice the pushed diagonal
(tangency multiplicity $2$) leaves
\[
20H^2+32HR+3R^2=(2H+3R)(10H+R),
\]
i.e.\ the residual cycle is a divisor of class $(10,1)$ on the
$(2,3)$ hypersurface --- the same answer by an independent route.
\end{remark}

\subsection{The explicit parametrization}\label{sec:explicit}

For computations (and for the external verification of
\S\ref{sec:artifacts}) it is convenient to have the composite map in
closed staged form.  With $L=\{W=0\}$, parametrize $L$ by
$y(t)=(t{:}1{:}0)$, and suppose (after the harmless linear specialization
described in \S\ref{sec:artifacts}, A4) that $e=(1{:}0{:}0)$ lies on every
conic $F(-,y(t))$.  Writing $A_{ij}(t)$ for the restricted coefficients of
the conic:
\begin{align*}
\textbf{(I)}\quad
&x^*(t,m)=\bigl(A_{11}+mA_{12}+m^2A_{22},\ -(A_{01}+mA_{02}),\
-m(A_{01}+mA_{02})\bigr)
\end{align*}
is the second intersection of the chord of direction $(0,1,m)$ through
$e$ --- a birational parametrization of $S$ by $(t,m)$.
\begin{align*}
\textbf{(II)}\quad
&f^*=F(x^*(t,m),-),\qquad
(\alpha,\beta,\gamma)=\nabla_y f^*\big|_{y(t)},\qquad
z^*=(\gamma,0,-\alpha),\\
&f^*\bigl(s\,y(t)+u\,z^*\bigr)=u^2\bigl(\mathsf A\,s+\mathsf B\,u\bigr),
\qquad
r(t,m)=\mathsf B\,y(t)-\mathsf A\,z^*
\end{align*}
($z^*$ is a second point of the tangent line, since
$\alpha\gamma+\gamma(-\alpha)=0$; the two leading coefficients vanish
identically, expressing tangency, and $r$ is the residual).
\begin{align*}
\textbf{(III)}\quad
&d=(0,1,n),\qquad Q=F(-,r),\qquad
x'(t,m,n)=Q(d)\,x^*-\bigl(Q(x^*{+}d)-Q(x^*)-Q(d)\bigr)\,d .
\end{align*}
Then $\Phi(t,m,n)=(x'(t,m,n);\,r(t,m))$ is a dominant rational map
$\mathbf A^3\dashrightarrow X$ of degree $20$.  Stage (III) is the chord
construction on the conic fiber through the marked point $x^*$, and its
correctness is the one-line \emph{reflection identity} for any ternary
quadratic form $Q$ with polarization $B_Q$:
\[
Q\bigl(Q(d)\,x-B_Q(x,d)\,d\bigr)=Q(d)^2\,Q(x),
\]
so $F(x',r)=Q(d)^2F(x^*,r)=0$.
The degree over $X$ equals $\deg(T/\PP^2_y)=T\cdot f_y=20$.

\begin{remark}[Scope]\label{rem:scope}
The construction uses the plane-cubic fibration $\rho$, which exists
precisely because the second degree is $3$.  It therefore says nothing
about arbitrary standard conic bundles with smooth degree-$9$ discriminant
(the matrix entries need not be all cubic), nor about bidegrees $(2,d)$,
$d\ge4$, where the fibers of $\rho$ have genus $>1$ and the tangent line
meets the fiber in $d-2>1$ residual points (a correspondence, not a map).
Harris's question for $d$ large \cite[p.~24]{Harris} remains open.
\end{remark}

\begin{remark}[Byproducts in the repository]
Two independent negative results certified in the same project delimit how
sharp the construction is: on a general $X$, \emph{every} member of
$|\Hx|$ is a $K3$ double plane (its branch sextic never acquires a point
of multiplicity $\ge3$), and \emph{no} member of $|\Hx+\Hy|$ has rational
normalization.  So the minimal-class searches fail provably, and the first
horizontal rational surface found lives in $|10\Hx+\Hy|$ --- a singular
member, as forced by adjunction (Lemma~\ref{lem:pic} ff.).
\end{remark}

\section{Verification artifacts}\label{sec:artifacts}

The proof above is complete as mathematics except where it invokes three
kinds of finite algebraic facts: the universal identity
(Theorem~\ref{thm:identity}), the primitivity of $q$, and the existence of
explicit witnesses (Lemma~\ref{lem:openness}).  These are machine-checked.
This section is a complete ledger.  All paths are relative to the
repository \repo; logs of every run are committed beside the scripts.

\subsection*{A1 --- the certificate of record}
\texttt{problems/B-conic-bundle-multisections/certificates/}
\texttt{tangent\_residual\_local\_checks.py} (with \texttt{.log}):
\begin{itemize}
\item verifies \eqref{eq:identity} \emph{fully symbolically} over the
universal ring $\ZZ[a_3,\dots,k]$ in ten independent coefficient symbols
(the $W^0,W^1$ coefficients of $\calR+\Delta f$ vanish identically; the
quotient by $W^2$ is linear);
\item verifies each coefficient of $q_f$ is homogeneous of degree $5$ and
that the three have $\gcd 1$ universally, and again after a recorded
substitution by ten quadratic forms (so the general $(2,3)$ equation has
primitive $q$);
\item supplies the smooth-fiber witness of Lemma~\ref{lem:openness}:
an explicit $(F_0,x_0)$ with the cubic smooth (unit Gr\"obner bases in
three charts), three residual points computed, verified to lie on the
cubic, on the tangent lines, on $\{q=0\}$, off $L$, pairwise distinct;
\item records the divisor arithmetic $(10,3)-2(0,1)=(10,1)$, degree $20$.
\end{itemize}
The full prose proof with all genericity qualifications is
\texttt{certificates/tangent\_residual\_theorem.md}.

\subsection*{A2 --- independent re-derivation (audit)}
During the audit of this session, the identity was re-derived from scratch
in a fresh SymPy session sharing no code with A1: the $W^0$ divisibility,
the cofactor being the \emph{classical} discriminant (both sign
conventions tested), degree-$5$ homogeneity, and universal $\gcd=1$ were
reconfirmed independently.
% (transcript: session audit, 2026-07-22)

\subsection*{A3 --- independent end-to-end geometric check (audit)}
On an explicit cubic with rational points on $L$
($g=(u-v)(u-2v)(u+v)$, generic $h,i,j,k$), the three tangent residuals
were computed \emph{from the raw geometry} (tangent lines; factoring the
double root out of the restricted cubic) and verified to be distinct, off
$L$, and to lie on $\{q_f=0\}$ --- confirming the geometric
interpretation, not merely the algebra.  A second check drew a random
$(2,3)$ form over $\QQ$ and verified the witness conditions of
Lemma~\ref{lem:openness} at $x_0=(1,2,3)$, plus $\gcd(q_U,q_V,q_W)=1$ by
restriction to two lines.

\subsection*{A4 --- external random-coefficient instance}
\texttt{problems/B-conic-bundle-multisections/external\_check/}: a full
end-to-end falsification test on coefficients drawn from
\texttt{/dev/urandom} (recorded), with four coefficients zeroed so that
the section $e=(1{:}0{:}0)$ exists over $\QQ$ and the composite map has
exact rational coefficients (a codimension-$4$ linear slice needed only
because Tsen's theorem is not effective over $\QQ$; over $\CC$ the theorem
needs no such conditioning).  The harness \texttt{build\_explicit\_map.py}
constructs the staged map of \S\ref{sec:explicit} and verifies, exactly:
\begin{center}
\begin{tabular}{@{}lll@{}}
\toprule
check & statement & method\\
\midrule
C1 & $F(x^*,y(t))\equiv0$ & symbolic\\
C2 & tangency (two leading coefficients vanish) & symbolic\\
C3 & $f^*(r)\equiv0$ & full $100\times53$ grid$^{\dagger}$\\
C4 & $q_{f^*}(r)\equiv0$ & full $62\times37$ grid$^{\dagger}$\\
C5 & reflection identity & generic symbols\\
C6 & $F(\Phi)=0$ at sample rational points & exact evaluation\\
C7 & Jacobian rank $3$ (dominance) & exact symbolic differentiation\\
C8 & fiber over random $y_0$: squarefree of degree $20$ & resultant\\
\bottomrule
\end{tabular}
\end{center}
{\small $^{\dagger}$Deterministic proofs: a polynomial of bidegree at most
$(D_1,D_2)$ vanishing at a $(D_1{+}1)\times(D_2{+}1)$ grid of distinct
rational nodes is zero; the bidegree bounds $(99,52)$ and $(61,36)$ are
derived in the harness.}\\[2pt]
All eight pass (\texttt{run\_log.txt}); the explicit map instance is
\texttt{map\_instance.txt}, and \texttt{README\_CHECK.md} gives
CAS-agnostic instructions for third-party verification with no trust in
the pipeline.

\subsection*{A5 --- unconditioned instance (no slice)}
\texttt{external\_check/unconditioned\_check.py} (with \texttt{.log})
addresses the natural objection that the codimension-$4$ slice in A4 might
matter: it runs the full mechanism on a fresh entropy draw with \emph{no}
coefficient zeroed.  With no rational point available, it works in the
cubic \'etale algebra $K=\QQ[\theta]/(g_0)$, $g_0=F(x_0,-)|_L$, where the
three marked points live; reduction modulo the full cubic verifies each
identity at all three conjugate points simultaneously.  It checks:
primitivity of $q_F$; the witness conditions of Lemma~\ref{lem:openness};
tangency, $F(x_0,r)=0$, $q(r)=0$, and $r\notin L$ (via
$\Res(g_0,r_W)\ne0$) in $K$; the stage-(III) point $F(x',r)=0$ in $K$; and
the squarefree degree-$20$ fiber.  All pass.  The slice in A4 is thereby
confirmed to play no role beyond supplying $\QQ$-rationality of the
composite parametrization (Tsen's theorem being non-effective over
$\QQ$); by $\mathrm{PGL}_3\times\mathrm{PGL}_3$-equivariance the sliced
family $\{X\supseteq\{e\}\times L\}$ is in no special position for
anything the proof uses.

\subsection*{Prose-only steps}
For completeness: the steps of the proof resting on standard mathematics
rather than computation are Tsen's theorem, Bertini/irreducibility of $S$,
Grothendieck--Lefschetz (Lemma~\ref{lem:pic}), the group-law remark, the
properness argument of Lemma~\ref{lem:openness}, and the Chow-theoretic
cross-check.  Each is classical.

\subsection*{Independence and audit trail}
A1 was rerun in a clean environment during the audit (Macaulay2 1.26.06 /
SymPy 1.14.0) with output matching the committed logs; A2--A3 share no
code with A1; A4 takes its input from an external entropy source.  The
literature verification (the boundary $\deg\Delta\le8$, the status of
\cite{Mella} as an unpublished preprint whose degree-$\le8$ consequence
has published backing via \cite{KollarMella,Prokhorov}, the method and
status of \cite{DCV}, and the absence of any prior statement of
Theorem~\ref{thm:main}) was performed against primary sources on
2026-07-22.

\begin{thebibliography}{99}

\bibitem{ArtinMumford}
M.~Artin, D.~Mumford,
\emph{Some elementary examples of unirational varieties which are not
rational}, Proc.\ London Math.\ Soc.\ (3) \textbf{25} (1972), 75--95.

\bibitem{Beauville}
A.~Beauville,
\emph{Vari\'et\'es de Prym et jacobiennes interm\'ediaires},
Ann.\ Sci.\ \'Ecole Norm.\ Sup.\ (4) \textbf{10} (1977), 309--391.

\bibitem{BvB}
C.~B\"ohning, H.-C.~Graf von Bothmer,
\emph{On stable rationality of some conic bundles and moduli spaces of
Prym curves}, Comment.\ Math.\ Helv.\ \textbf{93} (2018), 133--155.

\bibitem{CGM}
L.~Casarotti, S.~Gammelgaard, A.~Massarenti,
\emph{On the unirationality of conic bundles with discriminant of degree
eight}, preprint, \href{https://arxiv.org/abs/2511.17213}{arXiv:2511.17213}
(2025).

\bibitem{CTV}
J.-L.~Colliot-Th\'el\`ene, C.~Voisin,
\emph{Cohomologie non ramifi\'ee et conjecture de Hodge enti\`ere},
Duke Math.\ J.\ \textbf{161} (2012), 735--801.

\bibitem{DCV}
A.~Del Centina, L.~Verdi,
\emph{Su alcune variet\`a unirazionali di dimensione tre},
Atti Accad.\ Naz.\ Lincei Rend.\ Cl.\ Sci.\ Fis.\ Mat.\ Nat.\ (8)
\textbf{69} (1980), no.~6, 338--345.

\bibitem{Harris}
J.~Harris,
\emph{Rationality, unirationality and rational connectivity},
lecture notes, AMS Summer Institute in Algebraic Geometry (Seattle),
\url{https://www.math.columbia.edu/~thaddeus/seattle/harris1.pdf}.

\bibitem{KollarMella}
J.~Koll\'ar, M.~Mella,
\emph{Quadratic families of elliptic curves and unirationality of degree
1 conic bundles}, Amer.\ J.\ Math.\ \textbf{139} (2017), 915--936.

\bibitem{MassarentiMella}
A.~Massarenti, M.~Mella,
\emph{On the birational geometry of conic bundles over the projective
space}, Math.\ Nachr.\ \textbf{297} (2024), 1208--1220.

\bibitem{Mella}
M.~Mella,
\emph{On the unirationality of 3-fold conic bundles}, preprint,
\href{https://arxiv.org/abs/1403.7055}{arXiv:1403.7055} (2014).

\bibitem{Prokhorov}
Yu.~Prokhorov,
\emph{The rationality problem for conic bundles},
Russian Math.\ Surveys \textbf{73} (2018), no.~3, 375--456.

\end{thebibliography}

\end{document}
